\section{Extensions toward closed-loop agents}
\label{sec:extensions}

A closed-loop study would ask whether an agent can use rejection witnesses to repair a plan, obtain missing evidence, and complete useful work. A failed applicability check should identify the affected source and obligation. The fixed-action results establish neither the semantic quality of that feedback nor a model's ability to act on it.

A task-specific \emph{ContextPack} would carry selected source and proof dependencies, with local aliases mapped to original content identities, instead of whole collections. Semantic minification must preserve required obligations, source maps, unsupported facets, and current-root dependencies. Retrieval relevance cannot certify that an omitted constraint is irrelevant. Evaluation should measure context size, token use, compilation and checking overhead, invalidation, and authorized task utility. No compressor or token-saving result is reported here.

Further directions include revision-aware program states, exact or abstract transitions, learned next-action and inverse-trace proposals, and semantics-preserving refactoring with parallel validation and checked test reuse. Their evaluation must account for both preservation of the required guarantees and checking cost.
